\section{Key Stakeholders \& Blocs}

\lead{The following section summarises the different positions and initiatives taken across regions on the issue of child soldier use and recruitment, while providing an overview of some of the actions that have already been taken by international organizations, NGOs, and civil society.}

\subsection{Western Bloc}

Western states such as the United States, the United Kingdom, France, and Germany have tackled the issue primarily through legislative and financial leverage rather than through direct intervention. For instance, the United States enacted the Child Soldiers Prevention Act (CSPA) in 2008, requiring the publication of an annual list of governments that recruit child soldiers and prohibits certain types of U.S. Security assistance to such countries that engage in this practice\footnote{Weber, Michael A. "Child Soldiers Prevention Act: Security Assistance Restrictions." Congressional Research Service. October 11, 2023.\url{https://www.congress.gov/crs-product/IF10901}}. The law was later strengthened by the Child Soldier Prevention Act of 2018.\footnote{\emph{Ibid.}} In reality, however, the leverage has been political and inconsistently applied at best. The Biden Administration, for instance, had identified 17 governments complicit in child soldier use or recruitment, restricting arms sales and military assistance to nine of them while still granting partial restrictions to seven others, highlighting the relationship between Western commitments and strategic expectations made in respect of national interests.\footnote{Stimson Center. "Taking Stock of Biden's Delayed — but Encouraging — Progress on Child Soldiers Prevention." February 20, 2024.\url{https://www.stimson.org/2023/taking-stock-of-bidens-delayed-but-encouraging-progress-on-child-soldiers-prevention/}}

On the contrary, France has taken an operational role as the United Nations Security Council penholder on conflicts regarding the Central African Republic and the Democratic Republic of Congo, regularly advocating for the implementation of child protection provisions within UN peacekeeping mandates.\footnote{Watchlist on Children and Armed Conflict. "Children and Armed Conflict Monthly Update — June 2026." June 2026.\url{https://watchlist.org/publications/children-and-armed-conflict-monthly-update-june-2026/}}

The Western Bloc broadly supports the Vancouver Principles of 2017, which call on troop and police contributing countries to UN peacekeeping operations to strengthen child protection training and safeguards into their deployment.\footnote{Government of Canada. "Vancouver Principles on Peacekeeping and the Prevention of the Recruitment and Use of Child Soldiers." November 15, 2017.\url{https://www.international.gc.ca/world-monde/issues_development-enjeux_developpement/human_rights-droits_homme/principles-vancouver-principes.aspx?lang=eng}} This reflects the bloc’s preference for promoting child protection within current security measures rather than creating new standalone enforcement bodies. Western countries have also Security Council briefings to spotlight specific crises, with the United States flagging recruitment by both warring parties in Sudan’s conflict, where young boys have been recruited as child soldiers, alongside cases of child abduction, detention, and torture documented in Ukraine. \footnote{United Nations Security Council. "Facing Record-High Violations in 2023, Security Council Explores Ways to Bolster Norms to Protect Children in Armed Conflict." SC/15745. June 2024.\url{https://press.un.org/en/2024/sc15745.doc.htm}}The Western bloc generally uses strong normative language, sanctions-linked accountability for reported parties, and well-resourced reintegration programming.

\subsection{Latin American Bloc}

In Latin America, Colombia has been singled out as the country most affected by the issue at hand. Colombia has been experiencing a lot of internal conflict, and even though a peace agreement was signed in 2016 between the government and the nation’s most important guerrilla group, child soldier recruitment has increased by 300\% over the past five years, according to data from the year 2024, which was verified by the UN. This problem is especially pressing across poor rural areas in the country where children do not have a lot of educational opportunities or access to jobs\footnote{Colarullo Giacomo. \emph{Child recruitment and use by armed groups in Colombia quadruples over five years}. UNICEF. 2026. \url{https://www.unicef.org/press-releases/child-recruitment-and-use-armed-groups-in-colombia-quadruples-over-five-years}}. Colombia has tried to solve it primarily by implementing a DDR and child recruitment prevention model mainly centered around rehabilitating victims and expanding socio-economic opportunities for them to be less vulnerable to re-recruitment\footnote{International Organization for Migration. Colombia: DDR and Child Soldier Issues: A Monthly Review. Reliefweb. 2018. \url{https://reliefweb.int/report/colombia/colombia-ddr-and-child-soldier-issues-monthly-review-february-2018}}.

\subsection{Africa and the Middle East Bloc}

As touched upon in previous sections in this study guide, the West and Central Africa region stands among those most affected by human rights violations against children in war contexts. According to UNICEF, from the year 2016, more than 21,000 children have been confirmed by the UN as having been recruited and used by armed forces in West and Central Africa\footnote{Bisin Sandra. The West and Central Africa region among the most affected by grave violations against children in armed conflict. UNICEF. 2021. \url{https://www.unicef.org/press-releases/west-and-central-africa-region-among-most-affected-grave-violations-against-children}}. Within the broader African and Middle East region, the Democratic Republic of Congo, Somalia, Sudan, South Sudan, Nigeria, the Central African Republic, Syria, Yemen, Iraq, Israel and Palestine have all suffered from the problem of child soldiers in one way or the other.

As a result, it would make the most sense for many of the countries making up the African and Middle Eastern Bloc to act as strong advocates for action on the matter. They would mainly pay attention to stopping recruitment, guaranteeing rehabilitation, and securing international assistance, always bearing in mind the importance of safeguarding state sovereignty. On top of this, they would most likely acknowledge the root causes of child soldier use and recruitment and attempt to address them. A lot of these countries have undergone prolonged conflicts and have had to endure the presence of various non-state armed groups operating completely outside the government and international law. If this was not enough, there is also a lot of poverty in some of the regions affected by the child soldier problem, and this means that a lot of the children living here are particularly vulnerable to recruitment.

\subsection{Asian Bloc}

In Asia, countries such as Myanmar, the Philippines, or Afghanistan have experienced cases of child soldier use and recruitment. At the Asia-Pacific Conference on the Use of Children as Soldiers held in the year 2000 in Nepal, which culminated in the Kathmandu Declaration on the Use of Children as Soldiers, it was recognised that Asia is the second region most impacted by child soldier recruitment after Africa and that it is expected that approximately 75,000 children are involved in armed conflicts in some way or the other across countries such as the ones that were previously mentioned\footnote{Pillai R. V. Asia-Pacific Conference on the Use of Children as Soldiers. Nepal. 2000. \url{file:///C:/Users/Usuario/Downloads/Asia\%20Pacific\%20Conference\%20on\%20the\%20use\%20of\%20Children\%20as\%20Soldiers.pdf}}. Hence, similarly to Africa and the Middle East, the Asian countries should actively call for the implementation of more effective solutions to the issue of child soldiers.

\subsection{International Organizations}

Within the UN, UNICEF stands out as one of the most active agencies seeking to address the problem of child soldier use and recruitment. As mentioned in previous sections, it has been involved in negotiating child soldier releases as well as in providing psychological support to former child soldiers in an attempt to reintegrate them back into society. Simultaneously, the Office of the Special Representative of the Secretary-General for Children and Armed Conflict tries to encourage governments and armed groups to stop using and recruiting minors and is also in charge of monitoring violations against children’s rights in contexts of war\footnote{Office of the Special Representative of the Secretary-General for Children and Armed Conflict. \emph{Children Have Rights. Children Want Peace. Prove It Matters}. \url{https://childrenandarmedconflict.un.org/en}}.

At the regional level, the European Union has played an active role in protecting the rights of former child soldiers by funding various education and reintegration programs in countries like Uganda, the Democratic Republic of Congo, Colombia, Syria, and Sudan\footnote{European Union External Action. \emph{War is no Place for Children}. 2020. \url{https://www.eeas.europa.eu/eeas/war-no-place-children_en}}. More concretely, the EU has supported child soldiers in the Central African Republic through the provision of temporary care and reintegration kits as well as vocational training\footnote{European Union External Action. \emph{Child Soldiers: Children need Protection, not Trauma}. 2023. \url{https://www.eeas.europa.eu/eeas/child-soldiers-children-need-protection-not-trauma_en}}. Another EU initiative which has also demonstrated a strong commitment towards the rehabilitation of child soldiers is the decision to fund a project organised by UNICEF aiming to fully reintegrate children formerly associated with armed groups in South Sudan\footnote{Directorate-General for International Partnerships. \emph{Ex-post evaluation provision of reintegration services for children formerly associated with armed forces and groups in South Sudan}. European Commission. 2023. \url{https://international-partnerships.ec.europa.eu/publications-library/ex-post-evaluation-provision-reintegration-services-children-formerly-associated-armed-forces-and_en}}. Similarly, the African Union stands out as a key fighter against child soldier use and recruitment alongside the EU since it has led missions such as the African Union Transition Mission in Somalia (ATMIS), focused on training 30 military officers on how to protect children in war contexts and on implementing measures aimed at reducing the chances of children being recruited or abused during armed conflicts\footnote{ATMIS. \emph{ATMIS Troops Trained on Child Protection in Armed Conflict}. 2024. \url{https://atmis-au.org/atmis-troops-trained-on-child-protection-in-armed-conflict/}}. Finally, while organizations such as ASEAN, the OAE, or the Arab League have been less active on the matter, they have also positioned themselves against it.

\subsection{NGOs and Civil Society}

Some of the most active NGOs focused on tackling the issue of child soldier use and recruitment are Save the Children, War Child, Child Soldiers International, Human Rights Watch, and the International Committee of the Red Cross\footnote{Elder Shannon. \emph{Top Organizations Helping Child Soldiers Recover}. The Borgen Project. 2017. \url{https://borgenproject.org/organizations-helping-child-soldiers-recover/}}.

Save the Children, for instance, has been helping former child soldiers in countries such as the Democratic Republic of Congo through economic and psychological assistance\footnote{Garmirian Claire. \emph{Children Recruited into Armed Groups in DR Congo Describe Abduction, Drugging, and Violence as Call for Greater Action}. Save the Children. Kinshasa. 2026. \url{https://www.savethechildren.org/us/about-us/media-and-news/2026-press-releases/drc-children-recruited-into-armed-groups}}. Save the Children has been particularly important in leading advocacy efforts through initiatives such as the Stop the War on Children Report, which documents the atrocities suffered by children in war zones, emphasising how child soldier recruitment is a grave violation of children’s rights\footnote{Save the Children. \emph{Child Soldiers: The Tragic End of Childhood for Boys and Girls in Conflict}. 2026. \url{https://www.savethechildren.org/us/charity-stories/child-soldiers}}. In a similar manner, Child Soldiers International and Human Rights Watch have also stood out for publishing global reports documenting child soldier use and recruitment rates\footnote{Human Rights Watch. Child Soldiers Worldwide. 2012. \url{https://www.hrw.org/news/2012/03/12/child-soldiers-worldwide}}. Child Soldiers International, in particular, has been an active proponent of raising the global recruitment age to 18. Finally, Watchlist on Children and Armed Conflict is a less famous NGO but is also highly relevant since it has long focused on addressing the problem of child soldiers by assisting the UN Security Council Working Group on Children and Armed Conflict through research and recommendations\footnote{Watchlist on Children and Armed Conflict. \emph{About}. 2026. \url{https://watchlist.org/about/advocate/}}.
